\Chapter{Összegzés}

A jelen munka egy SQL:1999 kompatibilis lekérdezéseket EF Core 8.0 LINQ metódusszintaxisra fordító eszköz tervezését és megvalósítását mutatta be. A projekt két alapvető kérdést vizsgált: milyen SQL nyelvi elemek fordíthatók egyértelműen LINQ-ra, és milyen architektúra teszi lehetővé egy ilyen fordító megbízható és bővíthető megvalósítását.

Az első kérdésre a válasz a projekt hatókörén keresztül rajzolódott ki. A \texttt{SELECT} utasítás teljes egészében leképezhető: a \texttt{WHERE} feltételek, a \texttt{JOIN} műveletek, a \texttt{GROUP BY} aggregáció, a halmazműveletek, az allekérdezések és a \texttt{CASE WHEN} kifejezések mind rendelkeznek egyértelmű LINQ megfelelővel. A \texttt{DELETE}, \texttt{INSERT} és \texttt{UPDATE} utasítások szintén fordíthatók, bár az EF Core 8.0 bulk-műveletei (\texttt{ExecuteDeleteAsync}, \texttt{ExecuteUpdateAsync}) és a change tracker alapú műveletek (\texttt{Add}, \texttt{SaveChanges}) közötti különbség tudatos tervezési döntést igényelt. Bizonyos SQL konstrukciók viszont nem fordíthatók le megbízhatóan: a \texttt{FULL OUTER JOIN}, a közös táblakifejezések (CTE) és az ablakfüggvények nem rendelkeznek olyan LINQ megfelelővel, amelyet az EF Core minden provideren szerver oldalon SQL-re fordítana.

A második kérdésre a háromrétegű architektúra adott választ. Az \\ANTLR4-alapú parser, a saját LINQ AST és a rekurzív kódgenerátor szétválasztása lehetővé tette, hogy minden réteg önállóan fejleszthető és tesztelhető legyen. Az AST réteg különösen fontos szerepet játszott, mivel a parse fa és a LINQ AST szerkezetének különbsége nem pusztán technikai részlet, hanem a fordítás lényege. Ahol a parse fa az SQL szintaxis minden elemét tartalmazza, a LINQ AST a LINQ metódusláncot követi, és csak azokat az információkat hordozza, amelyek a kódgeneráláshoz szükségesek. A rétegek szétválasztása később azt is lehetővé tette, hogy a kódgenerálás a fordító többi részének módosítása nélkül cserélhető komponenssé váljon.

A fejlesztés során felmerülő kihívások visszatérő mintát mutattak, az első, kézenfekvőnek tűnő megközelítés szinte minden esetben valamilyen szélső esetben megbukott. A \texttt{LIKE} operátor regex-alapú fordítása kliensoldalon ragadt, az \texttt{IN (SELECT ...)} \texttt{Contains()}-alapú implementációja nullable típusoknál típusillesztési hibát dobott, a scope kezelés első változata lambda névütközéseket okozott mélyen egymásba ágyazott allekérdezéseknél. Ezek a tapasztalatok rávilágítottak arra, hogy egy fordító helyessége nem igazolható pusztán a happy path eseteken végzett ellenőrzéssel, a regressziók megelőzése érdekében a tesztelési infrastruktúra legalább annyira fontos, mint maga a fordítási logika.

A tesztelési megközelítés ennek megfelelően háromszintű lett. A string-szintű transpiler tesztek gyors visszajelzést adnak regressziókra. A Roslyn-alapú szemantikai tesztek igazolják, hogy a generált kód nem csupán szintaktikailag helyes, hanem valóban ugyanazt az eredményt adja, mint az eredeti SQL. A fuzz tesztek 5500 véletlenszerűen generált utasításon ellenőrzik a fordító robusztusságát olyan kombinációkban, amelyeket kézzel nehéz lenne megkonstruálni.

A projekt egy Blazor WebAssembly alapú webes demóval egészült ki, amely a fordítót interaktívan elérhetővé teszi a böngészőből, ANTLR parse fa és LINQ AST vizualizációval kiegészítve. Ez nemcsak demonstrációs célokat szolgál, hanem lehetővé teszi, hogy a fordító működése a parse fától a generált kódig vizuálisan is követhető legyen.

A kész fordító egy további lépésben keretrendszerré bővült. A kódgenerálás a köztes fa csomópontjaiból egy önálló, cserélhető komponensbe került, amelynek követelményeit az \\ \texttt{ICodeGenerator} interfész rögzíti. Ehhez a köztes fát is célnyelv-függetlenné kellett tenni, mivel az eredeti megoldásban a visitor több helyen már a fa felépítésekor C\# kódot állított elő. Az átalakítás után a C\# kódot a \texttt{CSharpCodeGenerator} osztály állítja elő, egy másik célnyelv támogatásához pedig elegendő egy újabb generátor megírása, a grammatika, a parser és a visitor változtatás nélkül újrafelhasználható. Az átalakítást a meglévő tesztkészlet tette biztonságossá, mivel a 267 teszteset elvárt kimenetén nem kellett módosítani. A választott megoldás a célnyelvek közötti szintaktikai különbségeket kezeli. A szemantikai eltérések, például a szerveroldali kiértékelés hiánya a Java Stream API esetében, egy konkrét második generátor megvalósításakor válna kezelhetővé, ami a munka folytatását jelenti.

A munka egyik tanulsága, hogy az SQL és a LINQ közötti leképezés a felszínen egyszerűbbnek tűnik, mint amilyen valójában. Mindkét nyelv deklaratív absztrakciókat kínál az adatok felett, de az absztrakciók mögötti modellek, a relációs algebra és az objektumgráf lényegesen különböznek egymástól. Egy fordítóprogram nem kerülheti meg ezt a különbséget, csupán szisztematikusan kezelheti. A projekt rámutatott, hogy ez a szisztematikus kezelés megvalósítható, de minden egyes konstrukciónál meg kell érteni, hogyan viselkedik a célplatform, az EF Core query provider, és milyen korlátokat szab a fordított kódra.
